\section{Experimental Protocol}
\label{sec:experiments}

\subsection{Evaluation suites}

\begin{itemize}
  \item \textbf{LIBERO in-distribution evaluation:} standard LIBERO suites will
  measure task competence before perturbation. Unperturbed LIBERO-Goal episodes
  will also provide the in-distribution reference for the corresponding
  LIBERO-PRO tasks.

  \item \textbf{LIBERO-PRO evaluation:} the official Goal-Language,
  Goal-Object, Goal-Swap and Goal-Task suites contain ten tasks with 50 initial
  states per task. Their matched evaluation comprises 2,000 episode identities
  per policy: four suites, ten tasks per suite and 50 initial states per task.

  \item \textbf{Project-specific Goal-Object-OOD evaluation:} a separate
  18-task suite contributes 50 initial states per task, for 900 matched episode
  identities per policy. We report it separately and do not present it as an
  official LIBERO-PRO suite.

  \item \textbf{IsaacLab evaluation:} the Franka environment will separate
  in-distribution trials from controlled changes in object identity, category,
  geometry, material, pose, lighting, background, camera and combinations of
  these factors. Asset-family-disjoint splits will support the unknown-object
  analysis.

  \item \textbf{Physical evaluation:} frozen simulation-selected checkpoints,
  monitors and thresholds will be transferred to the real Franka setup. Trial
  blocks will cover familiar objects, novel identities within known categories
  and held-out object categories.
\end{itemize}

Within each completed LIBERO protocol, an episode identity is the tuple of
suite, task, initial-state index and evaluation seed. Compared policies share
the complete identity set and instruction, and the campaigns use one evaluation
seed with 50 initial states per task. Observations are rendered at
\(480\!\times\!640\) and passed through each model's configured preprocessing.
The action budget is 250 environment actions; success is the environment task-
completion signal and exhausting the budget without completion is a failure.
Normal execution uses ten-action prefixes. The selected controller uses
\(H=10\), a five-action VLA prefix after mutual rejection and a 50-action
world-model eligibility latch. Branch thresholds are calibrated independently
before online evaluation and remain frozen on the matched test episodes.

\subsection{Baselines and comparison policy}

\begin{itemize}
  \item The matched LIBERO comparison comprises standalone HF SimVLA H10,
  standalone base world-model H10, fresh re-query U-VOWEL and the selected K1
  latch-50 U-VOWEL controller.

  \item Historical cached K1/K3 controllers and the CUDA legacy-cache
  reproduction are retained only as diagnostic ablations. Because they may
  retain rejected world-model state, they are excluded from the headline
  controller claim.

  \item Branch-specific ablations compare base and uncertainty-instrumented
  policies, monitor architectures and K1/K3 predictive-context variants under
  their corresponding matched manifests.

  \item Reproduced state-of-the-art policy baselines will use the same
  observations, prompts, action interface and episode manifests whenever their
  implementations permit. Published results obtained under different protocols
  will be reported in a separate comparison block and excluded from matched
  significance tests.

  \item The IsaacLab environment is project-specific. Its comparison will use
  strong policy families adapted to the same data and control interface rather
  than presenting an unmatched public-leaderboard claim.
\end{itemize}

\subsection{Metrics and statistical analysis}

\begin{itemize}
  \item Primary task metrics are episode success, per-task success and
  suite-level success. Paired outcome tables will count shared successes,
  shared failures and one-policy-only successes on identical episodes.

  \item Offline monitor metrics include AUROC, AUPRC, risk--coverage behavior,
  failure detection, success false alarms, detection time and warning lead time
  where the labels support it. A success false alarm is any alarm in a successful
  episode; failure detection is the fraction of failed episodes containing an
  alarm. Det@10, Det@25 and Det@50 measure failures first detected within 10\%,
  25\% and 50\% of the evaluated horizon. Never is the fraction of failed
  episodes with no alarm.

  \item Controller metrics include branch-query counts, escalation rate,
  world-model acceptance, mutual rejection, useful and harmful outcome
  reversals and latch activity.

  \item Efficiency metrics include branch inference time, monitor overhead,
  all-episode wall time, successful-episode wall time, throughput and peak GPU
  memory. Success-conditioned time will always be reported beside the
  all-episode statistic. Episode timing is runner wall time and excludes
  one-time model loading.

  \item The completed LIBERO campaigns contain one evaluation seed and repeated
  initial states nested within tasks. Exact McNemar tests describe paired
  episode differences but do not account for task clustering or model-selection
  history. We therefore report task-level results and do not treat episodes as
  independent task samples.
\end{itemize}

\subsection{Sim-to-real protocol and safety}

\begin{itemize}
  \item Simulation and hardware interfaces will be checked for camera geometry,
  color conversion, state representation, control frequency, workspace frame
  and action scaling.

  \item Policy order will be randomized within object and condition blocks. The
  VLA, world-model and fused controller will use the same trial manifest.
  Checkpoints, prompts, thresholds and stopping rules will remain frozen during
  the reported evaluation.

  \item Outcomes will distinguish success, drop, collision, timeout, automatic
  stop and human intervention. Independent workspace, velocity, force,
  watchdog and emergency-stop limits remain outside the learned controller.
\end{itemize}
